\section{Feature-Structure Disentanglement Framework}
\label{sec:fsd}

We introduce the \textbf{Feature-Structure Disentanglement (FSD)} framework, which provides a unified perspective on graph attention design by formally separating two orthogonal sources of node similarity.

\subsection{Motivation: The Hidden Assumption in GNNs}

Standard Graph Neural Networks operate on a fundamental assumption:

\begin{quote}
\textit{``Structural neighbors are the optimal source for information aggregation.''}
\end{quote}

This assumption underlies all message-passing GNNs, from GCN~\cite{kipf2017gcn} to GAT~\cite{velivckovic2018gat}:
\begin{equation}
\mathbf{h}_i^{(l+1)} = \text{UPDATE}\left(\mathbf{h}_i^{(l)}, \text{AGGREGATE}\left(\{\mathbf{h}_j^{(l)} : j \in \mathcal{N}_{\text{struct}}(i)\}\right)\right)
\end{equation}

However, this assumption fails in several important scenarios:

\textbf{Scenario 1: Fraud Detection with Camouflage.} Fraudsters deliberately connect to legitimate users to appear trustworthy. Structural neighbors include adversarial connections designed to \textit{mislead} aggregation.

\textbf{Scenario 2: High-Degree Hub Nodes.} In financial networks, large institutions connect to diverse entities. Their structural neighbors span multiple industries, risk profiles, and transaction patterns---aggregating from all neighbors dilutes node-specific signals.

\textbf{Scenario 3: Feature-Dominant Tasks.} When node labels primarily depend on intrinsic features (e.g., transaction amounts, account age) rather than network position, structural aggregation may introduce noise from feature-dissimilar neighbors.

These scenarios share a common pattern: \textbf{structural proximity does not imply feature similarity or label agreement}.

\subsection{Two Notions of Similarity}

We formalize the distinction between structure-based and feature-based similarity.

\begin{definition}[Structural Similarity]
\label{def:struct_sim}
Given graph $\mathcal{G} = (\mathcal{V}, \mathcal{E})$, the structural similarity between nodes $i$ and $j$ is:
\begin{equation}
S_{\text{struct}}(i, j) = f_{\text{struct}}(\mathcal{A}, i, j)
\end{equation}
where $f_{\text{struct}}$ can be:
\begin{itemize}
    \item \textbf{Adjacency}: $S_{\text{struct}}(i,j) = \mathcal{A}_{ij}$
    \item \textbf{Common neighbors}: $S_{\text{struct}}(i,j) = |\mathcal{N}(i) \cap \mathcal{N}(j)|$
    \item \textbf{Personalized PageRank}: $S_{\text{struct}}(i,j) = \text{PPR}(i,j)$
\end{itemize}
\end{definition}

\begin{definition}[Feature Similarity]
\label{def:feat_sim}
Given node features $\mathbf{X} \in \mathbb{R}^{N \times d}$, the feature similarity between nodes $i$ and $j$ is:
\begin{equation}
S_{\text{feat}}(i, j) = f_{\text{feat}}(\mathbf{x}_i, \mathbf{x}_j)
\end{equation}
where $f_{\text{feat}}$ can be:
\begin{itemize}
    \item \textbf{Cosine}: $S_{\text{feat}}(i,j) = \frac{\mathbf{x}_i^\top \mathbf{x}_j}{\|\mathbf{x}_i\| \|\mathbf{x}_j\|}$
    \item \textbf{Gaussian RBF}: $S_{\text{feat}}(i,j) = \exp(-\|\mathbf{x}_i - \mathbf{x}_j\|^2 / \sigma^2)$
    \item \textbf{Numerical (NAA)}: $S_{\text{feat}}(i,j) = -\frac{1}{d}\sum_{k=1}^{d} \frac{|\tilde{x}_{ik} - \tilde{x}_{jk}|}{\max(|\tilde{x}_{ik}|, |\tilde{x}_{jk}|, \epsilon)}$
\end{itemize}
\end{definition}

\textbf{Key Observation}: $S_{\text{struct}}$ and $S_{\text{feat}}$ are generally \textit{not aligned}. A node's structural neighbors may have very different features, and feature-similar nodes may be distant in the graph.

\subsection{Feature-Structure Alignment}

We quantify the relationship between structural and feature similarity.

\begin{definition}[Feature-Structure Alignment Score]
\label{def:fs_alignment}
The Feature-Structure Alignment Score $\rho_{\text{FS}}$ measures how well structural proximity correlates with feature similarity:
\begin{equation}
\rho_{\text{FS}} = \mathbb{E}_{(i,j) \in \mathcal{E}}[S_{\text{feat}}(i, j)] - \mathbb{E}_{(i,j) \notin \mathcal{E}}[S_{\text{feat}}(i, j)]
\end{equation}
\end{definition}

Intuitively:
\begin{itemize}
    \item $\rho_{\text{FS}} > 0$: Structural neighbors are more feature-similar than non-neighbors (\textbf{aligned})
    \item $\rho_{\text{FS}} \approx 0$: Structure and features are \textbf{independent}
    \item $\rho_{\text{FS}} < 0$: Structural neighbors are \textit{less} feature-similar (\textbf{misaligned})
\end{itemize}

\subsubsection{Computing $\rho_{\text{FS}}$ in Practice}

Computing the exact expectation over all non-edges is intractable for large graphs ($O(N^2)$ pairs). We provide an efficient sampling-based algorithm:

\begin{algorithm}[h]
\caption{Feature-Structure Alignment Score Computation}
\label{alg:fs_alignment}
\begin{algorithmic}[1]
\REQUIRE Graph $\mathcal{G} = (\mathcal{V}, \mathcal{E})$, features $\mathbf{X}$, sample size $M$, similarity function $S_{\text{feat}}$
\ENSURE $\hat{\rho}_{\text{FS}}$ (estimated alignment score)
\STATE \textbf{// Step 1: Compute edge similarity}
\STATE $\mathcal{S}_{\text{edge}} \leftarrow \emptyset$
\FOR{each edge $(i, j) \in \mathcal{E}$}
    \STATE $\mathcal{S}_{\text{edge}} \leftarrow \mathcal{S}_{\text{edge}} \cup \{S_{\text{feat}}(\mathbf{x}_i, \mathbf{x}_j)\}$
\ENDFOR
\STATE $\bar{S}_{\text{edge}} \leftarrow \text{mean}(\mathcal{S}_{\text{edge}})$
\STATE \textbf{// Step 2: Sample non-edges and compute similarity}
\STATE $\mathcal{S}_{\text{non-edge}} \leftarrow \emptyset$
\STATE $k \leftarrow 0$
\WHILE{$k < M$}
    \STATE Sample $i, j \sim \text{Uniform}(\mathcal{V})$ with $i \neq j$
    \IF{$(i, j) \notin \mathcal{E}$}
        \STATE $\mathcal{S}_{\text{non-edge}} \leftarrow \mathcal{S}_{\text{non-edge}} \cup \{S_{\text{feat}}(\mathbf{x}_i, \mathbf{x}_j)\}$
        \STATE $k \leftarrow k + 1$
    \ENDIF
\ENDWHILE
\STATE $\bar{S}_{\text{non-edge}} \leftarrow \text{mean}(\mathcal{S}_{\text{non-edge}})$
\STATE \textbf{// Step 3: Compute alignment score}
\STATE $\hat{\rho}_{\text{FS}} \leftarrow \bar{S}_{\text{edge}} - \bar{S}_{\text{non-edge}}$
\RETURN $\hat{\rho}_{\text{FS}}$
\end{algorithmic}
\end{algorithm}

\textbf{Complexity Analysis}: Step 1 requires $O(|\mathcal{E}| \cdot d)$ for computing all edge similarities. Step 2 requires $O(M \cdot d)$ for sampling and computing non-edge similarities. For sparse graphs where $|\mathcal{E}| \ll N^2$, the rejection sampling in Step 2 is efficient (acceptance rate $\approx 1$). Total complexity: $O((|\mathcal{E}| + M) \cdot d)$.

\textbf{Choice of $M$}: We use $M = \min(|\mathcal{E}|, 100000)$ to balance accuracy and efficiency. Empirically, $\hat{\rho}_{\text{FS}}$ stabilizes within $\pm 0.01$ after $M > 50000$ samples.

\textbf{Choice of $S_{\text{feat}}$}: We use cosine similarity as default: $S_{\text{feat}}(i,j) = \frac{\mathbf{x}_i^\top \mathbf{x}_j}{\|\mathbf{x}_i\| \|\mathbf{x}_j\|}$. For financial features with multi-scale values, we first apply log-normalization before computing cosine similarity.

\begin{theorem}[Estimation Error Bound for $\hat{\rho}_{\text{FS}}$]
\label{thm:estimation_error}
Let $\hat{\rho}_{\text{FS}}$ be the estimate from Algorithm~\ref{alg:fs_alignment} with $M$ non-edge samples. Assume $S_{\text{feat}}(i,j) \in [-1, 1]$ (e.g., cosine similarity). Then with probability at least $1 - \delta$:
\begin{equation}
|\hat{\rho}_{\text{FS}} - \rho_{\text{FS}}| \leq \sqrt{\frac{2\log(4/\delta)}{|\mathcal{E}|}} + \sqrt{\frac{2\log(4/\delta)}{M}}
\end{equation}

\textbf{Refined Bound for Large Graphs}: For sparse graphs with $|\mathcal{E}| = O(N)$ and $M = \Omega(N)$, the error bound simplifies to:
\begin{equation}
|\hat{\rho}_{\text{FS}} - \rho_{\text{FS}}| = O\left(\sqrt{\frac{\log(1/\delta)}{N}}\right)
\end{equation}

\textbf{Sample Complexity}: To achieve estimation error $\leq \epsilon$ with probability $\geq 1-\delta$, we need:
\begin{equation}
M \geq \frac{8\log(4/\delta)}{\epsilon^2}
\end{equation}
For $\epsilon = 0.01$ and $\delta = 0.05$, this gives $M \geq 75,000$, which justifies our default choice of $M = 100,000$.
\end{theorem}

\begin{proof}
The estimate $\hat{\rho}_{\text{FS}} = \bar{S}_{\text{edge}} - \bar{S}_{\text{non-edge}}$ consists of two independent sample means.

For $\bar{S}_{\text{edge}}$: This is the average over all $|\mathcal{E}|$ edges. By Hoeffding's inequality for bounded random variables in $[-1, 1]$:
\begin{equation}
\Pr\left[|\bar{S}_{\text{edge}} - \mathbb{E}[S_{\text{edge}}]| > \epsilon_1\right] \leq 2\exp\left(-\frac{|\mathcal{E}| \epsilon_1^2}{2}\right)
\end{equation}

For $\bar{S}_{\text{non-edge}}$: This is the average over $M$ uniformly sampled non-edges. Similarly:
\begin{equation}
\Pr\left[|\bar{S}_{\text{non-edge}} - \mathbb{E}[S_{\text{non-edge}}]| > \epsilon_2\right] \leq 2\exp\left(-\frac{M \epsilon_2^2}{2}\right)
\end{equation}

Setting $\epsilon_1 = \sqrt{2\log(4/\delta)/|\mathcal{E}|}$ and $\epsilon_2 = \sqrt{2\log(4/\delta)/M}$, and applying union bound:
\begin{equation}
\Pr\left[|\hat{\rho}_{\text{FS}} - \rho_{\text{FS}}| > \epsilon_1 + \epsilon_2\right] \leq \delta
\end{equation}

The refined bound for large graphs follows by noting that when $|\mathcal{E}| = O(N)$ and $M = c \cdot N$ for some constant $c \geq 1$, both terms scale as $O(\sqrt{\log(1/\delta)/N})$, giving the simplified asymptotic form.

The sample complexity bound is obtained by solving $\epsilon_2 \leq \epsilon/2$ for $M$, assuming $|\mathcal{E}|$ is sufficiently large that $\epsilon_1 \leq \epsilon/2$ also holds.
\end{proof}

\begin{remark}[Relationship to Homophily]
Feature-Structure Alignment ($\rho_{\text{FS}}$) is related to but distinct from graph homophily ($h$). Homophily measures \textit{label} similarity among neighbors:
\begin{equation}
h = \frac{|\{(i,j) \in \mathcal{E} : y_i = y_j\}|}{|\mathcal{E}|}
\end{equation}
While $\rho_{\text{FS}}$ measures \textit{feature} similarity among neighbors. When features are predictive of labels, high $\rho_{\text{FS}}$ often correlates with high $h$, but they capture different aspects of graph structure.
\end{remark}

\subsection{Aggregation Dilution: The Missing Factor}
\label{sec:aggregation_dilution}

While $\rho_{\text{FS}}$ captures the \textit{quality} of neighborhood information, it does not account for the \textit{quantity} of neighbors. In graphs with highly skewed degree distributions, high-degree nodes aggregate information from many neighbors, potentially \textbf{diluting} node-specific signals even when individual neighbor similarity is reasonable.

\begin{definition}[Aggregation Dilution]
\label{def:agg_dilution}
The Aggregation Dilution score $\delta_{\text{agg}}$ measures the degree-weighted information loss during mean aggregation:
\begin{equation}
\delta_{\text{agg}} = \mathbb{E}_{i \in \mathcal{V}}\left[ d_i \cdot \left(1 - S_{\text{agg}}(i)\right) \right]
\end{equation}
where $d_i = |\mathcal{N}(i)|$ is the degree of node $i$, and $S_{\text{agg}}(i)$ is the similarity between node $i$ and its \textit{aggregated} neighborhood:
\begin{equation}
S_{\text{agg}}(i) = \cos\left(\mathbf{x}_i, \frac{\bar{\mathbf{x}}_{\mathcal{N}(i)}}{\|\bar{\mathbf{x}}_{\mathcal{N}(i)}\|}\right), \quad \bar{\mathbf{x}}_{\mathcal{N}(i)} = \frac{1}{d_i}\sum_{j \in \mathcal{N}(i)} \mathbf{x}_j
\end{equation}

\textbf{Important Note}: We use similarity to the \textit{aggregated} neighbors rather than the average of pairwise similarities. This directly captures the effect of mean aggregation in GNNs: even if individual neighbors are similar, their mean may differ significantly from the ego node. Mathematically:
\begin{equation}
\mathbb{E}[d \cdot (1 - S_{\text{agg}})] \neq \mathbb{E}[d] \cdot (1 - \mathbb{E}[S_{\text{agg}}])
\end{equation}
This inequality arises because high-degree nodes tend to have more diverse neighbors (lower $S_{\text{agg}}$), causing the product to exceed the simple estimate.

\textbf{Estimation Error}: When computing $\delta_{\text{agg}}$ empirically from a sample of $n$ nodes, the estimation error satisfies with probability $1 - \delta$:
\begin{equation}
\left|\hat{\delta}_{\text{agg}} - \delta_{\text{agg}}\right| \leq d_{\max} \cdot \sqrt{\frac{2\log(2/\delta)}{n}}
\end{equation}
where $d_{\max}$ is the maximum node degree. For graphs with bounded degree $d_{\max} = O(\log N)$, the error is $O(\sqrt{\log N \cdot \log(1/\delta)/n})$.
\end{definition}

\textbf{Intuition}: $\delta_{\text{agg}}$ captures two compounding effects:
\begin{enumerate}
    \item \textbf{Dissimilarity to aggregation}: $(1 - S_{\text{feat}}(\mathbf{x}_i, \bar{\mathbf{x}}_{\mathcal{N}(i)}))$ measures how different a node is from its aggregated neighborhood
    \item \textbf{Degree amplification}: Multiplying by $d_i$ reflects that high-degree nodes suffer more from mean aggregation---their signal is averaged over many (potentially diverse) neighbors
\end{enumerate}

\begin{proposition}[Degree-Dependent Dilution]
\label{prop:degree_dilution}
For graphs with heterogeneous degree distributions, the aggregation dilution increases superlinearly with node degree:
\begin{equation}
\mathbb{E}[\delta_{\text{agg}} | d_i = d] \approx d \cdot (1 - \bar{S}_d)
\end{equation}
where $\bar{S}_d$ is the average similarity to aggregated neighbors for nodes of degree $d$. When neighbor features are drawn from diverse subpopulations (common in financial networks), $\bar{S}_d$ decreases as $d$ increases, causing $\delta_{\text{agg}}$ to grow faster than linearly.
\end{proposition}

\textbf{Implications for Method Selection}: High $\delta_{\text{agg}}$ indicates that mean-aggregation methods (GCN, GAT, NAA) will suffer significant information loss on high-degree nodes. In such cases:
\begin{itemize}
    \item \textbf{Sampling-based methods} (GraphSAGE): Aggregate from a fixed-size sample, bounding dilution
    \item \textbf{Concatenation methods} (H2GCN, MixHop): Preserve ego representation separately from aggregated neighbors
\end{itemize}

\begin{remark}[Extended FSD Framework]
The complete FSD diagnostic framework now includes three complementary metrics:
\begin{enumerate}
    \item $\rho_{\text{FS}}$ (Feature-Structure Alignment): Do structural neighbors have similar features?
    \item $h$ (Homophily): Do structural neighbors have the same label?
    \item $\delta_{\text{agg}}$ (Aggregation Dilution): Does mean aggregation cause information loss?
\end{enumerate}
Together, these metrics provide a comprehensive characterization for method selection.
\end{remark}

\subsection{Theoretical Analysis of Aggregation Dilution}
\label{sec:dilution_theory}

We now establish rigorous theoretical connections between $\delta_{\text{agg}}$ and GNN behavior, providing the foundation for principled method selection.

\subsubsection{Single-Layer Information Loss}

\begin{theorem}[Single-Layer Ego Deviation Bound]
\label{thm:single_layer}
Consider a single GCN layer with normalized adjacency $\tilde{\mathbf{A}}$ where $\tilde{a}_{ij} = 1/(d_i+1)$ for $j \in \mathcal{N}(i) \cup \{i\}$. Let $\mathbf{h}_i^{(0)} = \mathbf{x}_i$ be the initial feature and $\mathbf{h}_i^{(1)} = \sum_j \tilde{a}_{ij} \mathbf{h}_j^{(0)} \mathbf{W}$ be the representation after one layer. Then the ego deviation satisfies:
\begin{equation}
\left\| \mathbf{h}_i^{(1)} - \frac{1}{d_i+1}\mathbf{h}_i^{(0)}\mathbf{W} \right\| = \frac{d_i}{d_i+1} \left\| \bar{\mathbf{h}}_{\mathcal{N}(i)}^{(0)} \mathbf{W} \right\|
\end{equation}
where $\bar{\mathbf{h}}_{\mathcal{N}(i)}^{(0)} = \frac{1}{d_i}\sum_{j \in \mathcal{N}(i)} \mathbf{h}_j^{(0)}$ is the mean neighbor representation.

Furthermore, the deviation from ego is bounded by:
\begin{equation}
\left\| \mathbf{h}_i^{(1)} - \mathbf{h}_i^{(0)}\mathbf{W} \right\| \geq \frac{d_i}{d_i+1} \left\| \bar{\mathbf{h}}_{\mathcal{N}(i)}^{(0)} - \mathbf{h}_i^{(0)} \right\| \cdot \sigma_{\min}(\mathbf{W})
\end{equation}
where $\sigma_{\min}(\mathbf{W})$ is the minimum singular value of $\mathbf{W}$.
\end{theorem}

\begin{proof}
The GCN aggregation can be decomposed as:
\begin{align}
\mathbf{h}_i^{(1)} &= \sum_{j \in \mathcal{N}(i) \cup \{i\}} \tilde{a}_{ij} \mathbf{h}_j^{(0)} \mathbf{W} \\
&= \frac{1}{d_i+1} \mathbf{h}_i^{(0)} \mathbf{W} + \frac{1}{d_i+1} \sum_{j \in \mathcal{N}(i)} \mathbf{h}_j^{(0)} \mathbf{W} \\
&= \frac{1}{d_i+1} \mathbf{h}_i^{(0)} \mathbf{W} + \frac{d_i}{d_i+1} \bar{\mathbf{h}}_{\mathcal{N}(i)}^{(0)} \mathbf{W}
\end{align}

The deviation from the scaled ego representation is:
\begin{equation}
\mathbf{h}_i^{(1)} - \frac{1}{d_i+1}\mathbf{h}_i^{(0)}\mathbf{W} = \frac{d_i}{d_i+1} \bar{\mathbf{h}}_{\mathcal{N}(i)}^{(0)} \mathbf{W}
\end{equation}

Taking norms yields the first result. For the second inequality:
\begin{align}
\left\| \mathbf{h}_i^{(1)} - \mathbf{h}_i^{(0)}\mathbf{W} \right\| &= \left\| \frac{d_i}{d_i+1}(\bar{\mathbf{h}}_{\mathcal{N}(i)}^{(0)} - \mathbf{h}_i^{(0)})\mathbf{W} + \frac{-d_i}{d_i+1}\mathbf{h}_i^{(0)}\mathbf{W} \right\| \\
&\geq \frac{d_i}{d_i+1} \left\| (\bar{\mathbf{h}}_{\mathcal{N}(i)}^{(0)} - \mathbf{h}_i^{(0)})\mathbf{W} \right\| - \frac{d_i}{d_i+1}\left\|\mathbf{h}_i^{(0)}\mathbf{W}\right\|
\end{align}

Using $\|\mathbf{v}\mathbf{W}\| \geq \|\mathbf{v}\| \cdot \sigma_{\min}(\mathbf{W})$ completes the proof.
\end{proof}

\begin{corollary}[High-Degree Dominance]
\label{cor:high_degree}
As $d_i \to \infty$, the ratio $\frac{d_i}{d_i+1} \to 1$, implying:
\begin{equation}
\mathbf{h}_i^{(1)} \to \bar{\mathbf{h}}_{\mathcal{N}(i)}^{(0)} \mathbf{W}
\end{equation}
That is, high-degree nodes lose their ego information entirely after just one layer of aggregation.
\end{corollary}

\textbf{Connection to $\delta_{\text{agg}}$}: Using cosine similarity, $1 - S_{\cos}(\mathbf{x}_i, \bar{\mathbf{x}}_{\mathcal{N}(i)}) \propto \|\bar{\mathbf{x}}_{\mathcal{N}(i)} - \mathbf{x}_i\|^2$ for normalized features. Thus, $\delta_{\text{agg}}(i) = d_i \cdot (1 - S_{\cos})$ directly quantifies the deviation bound in Theorem~\ref{thm:single_layer}.

\subsubsection{Multi-Layer Over-smoothing Acceleration}

\begin{theorem}[Cumulative Deviation Bound]
\label{thm:multi_layer}
For an $L$-layer GCN with weight matrices $\{\mathbf{W}^{(l)}\}_{l=1}^L$ and ReLU activations, the cumulative deviation from the initial representation satisfies:
\begin{equation}
\left\| \mathbf{h}_i^{(L)} - \mathbf{h}_i^{(0)} \prod_{l=1}^L \mathbf{W}^{(l)} \right\| \leq C \cdot L \cdot \frac{d_i}{d_i+1} \cdot \sqrt{\delta_{\text{agg}}(i)}
\end{equation}
where $C = \prod_{l=1}^L \|\mathbf{W}^{(l)}\|$ is a constant depending on weight norms.
\end{theorem}

\begin{proof}[Proof Sketch]
We proceed by induction on $L$. The base case $L=1$ follows from Theorem~\ref{thm:single_layer}. For the inductive step, we bound the deviation at layer $l+1$ by the deviation at layer $l$ plus the new aggregation error. The key observation is that each layer contributes an additive term proportional to $\frac{d_i}{d_i+1} \cdot \|\bar{\mathbf{h}}_{\mathcal{N}(i)}^{(l)} - \mathbf{h}_i^{(l)}\|$. Summing over $L$ layers and using the relationship between Euclidean distance and cosine similarity yields the stated bound.
\end{proof}

\begin{remark}[Over-smoothing Acceleration]
Theorem~\ref{thm:multi_layer} reveals that nodes with high $\delta_{\text{agg}}$ over-smooth faster. Specifically:
\begin{itemize}
    \item High-degree nodes ($d_i \gg 1$): The factor $\frac{d_i}{d_i+1} \approx 1$ maximizes deviation
    \item Low similarity to neighbors: Large $(1 - S_{\cos})$ increases $\sqrt{\delta_{\text{agg}}}$
    \item The combination $d_i \cdot (1 - S_{\cos})$ captures both effects multiplicatively
\end{itemize}
This explains why over-smoothing is particularly severe in dense financial graphs with heterogeneous node features.
\end{remark}

\subsubsection{Sampling Bounds Dilution}

\begin{theorem}[GraphSAGE Dilution Bound]
\label{thm:sampling}
Consider GraphSAGE with uniform sampling of $k$ neighbors per node. For node $i$ with degree $d_i \geq k$, the effective aggregation dilution satisfies:
\begin{equation}
\delta_{\text{agg}}^{\text{SAGE}}(i) \leq k \cdot \max_{j \in \mathcal{N}(i)} \left(1 - S_{\text{feat}}(\mathbf{x}_i, \mathbf{x}_j)\right)
\end{equation}

In expectation over random samples:
\begin{equation}
\mathbb{E}\left[\delta_{\text{agg}}^{\text{SAGE}}(i)\right] = k \cdot \left(1 - \mathbb{E}_{j \sim \text{Unif}(\mathcal{N}(i))}[S_{\text{feat}}(\mathbf{x}_i, \mathbf{x}_j)]\right)
\end{equation}
\end{theorem}

\begin{proof}
GraphSAGE aggregates from a sampled subset $S \subseteq \mathcal{N}(i)$ with $|S| = k$:
\begin{equation}
\bar{\mathbf{x}}_S = \frac{1}{k} \sum_{j \in S} \mathbf{x}_j
\end{equation}

The effective dilution uses the sampled mean:
\begin{equation}
\delta_{\text{agg}}^{\text{SAGE}}(i) = k \cdot \left(1 - S_{\text{feat}}(\mathbf{x}_i, \bar{\mathbf{x}}_S)\right)
\end{equation}

By triangle inequality and properties of cosine similarity:
\begin{equation}
1 - S_{\text{feat}}(\mathbf{x}_i, \bar{\mathbf{x}}_S) \leq \max_{j \in S} \left(1 - S_{\text{feat}}(\mathbf{x}_i, \mathbf{x}_j)\right) \leq \max_{j \in \mathcal{N}(i)} \left(1 - S_{\text{feat}}(\mathbf{x}_i, \mathbf{x}_j)\right)
\end{equation}

The expectation follows from linearity of expectation over uniform sampling.
\end{proof}

\begin{corollary}[Sampling Decouples Degree]
\label{cor:sampling}
Crucially, $\delta_{\text{agg}}^{\text{SAGE}}(i)$ depends on the sampling budget $k$, \textbf{not} the original degree $d_i$. For a fixed $k$:
\begin{equation}
\delta_{\text{agg}}^{\text{SAGE}}(i) = O(k) \quad \text{regardless of } d_i
\end{equation}

This explains why GraphSAGE outperforms GCN on high-degree graphs like IEEE-CIS: sampling bounds dilution even when $d_i \gg k$.
\end{corollary}

\subsubsection{Concatenation Preserves Ego Information}

\begin{theorem}[Concatenation Eliminates Ego Dilution]
\label{thm:concat}
Consider H2GCN-style concatenation:
\begin{equation}
\mathbf{h}_i^{(l+1)} = \sigma\left(\left[\mathbf{h}_i^{(l)} \| \text{AGG}(\mathcal{N}(i)) \| \text{AGG}(\mathcal{N}^2(i))\right] \mathbf{W}\right)
\end{equation}

The ego component $\mathbf{h}_i^{(l)}$ is preserved in a dedicated subspace, with effective dilution:
\begin{equation}
\delta_{\text{agg}}^{\text{ego}}(i) = 0
\end{equation}

The aggregated components may still suffer dilution, but the ego representation remains intact.
\end{theorem}

\begin{proof}
In concatenation-based methods, the output can be written as:
\begin{equation}
\mathbf{h}_i^{(l+1)} = \sigma\left(\mathbf{h}_i^{(l)} \mathbf{W}_{\text{ego}} + \bar{\mathbf{h}}_{\mathcal{N}(i)}^{(l)} \mathbf{W}_{\text{neigh}} + \bar{\mathbf{h}}_{\mathcal{N}^2(i)}^{(l)} \mathbf{W}_{\text{2hop}}\right)
\end{equation}

where $\mathbf{W} = [\mathbf{W}_{\text{ego}}; \mathbf{W}_{\text{neigh}}; \mathbf{W}_{\text{2hop}}]$ is block-structured.

The gradient with respect to ego:
\begin{equation}
\frac{\partial \mathbf{h}_i^{(l+1)}}{\partial \mathbf{h}_i^{(l)}} = \text{diag}(\sigma') \cdot \mathbf{W}_{\text{ego}}
\end{equation}

This is independent of $d_i$ and does not involve any mean aggregation over neighbors. The ego information flows through a dedicated pathway, unaffected by neighborhood size.
\end{proof}

\begin{remark}[Why H2GCN Works on IEEE-CIS]
IEEE-CIS has high average degree (47.65) and high degree variance (CV=1.15). By Theorem~\ref{thm:concat}, H2GCN's concatenation strategy preserves ego information regardless of degree, while GCN/GAT/NAA suffer from dilution (Theorem~\ref{thm:single_layer}). This theoretical prediction matches our experimental observation that H2GCN outperforms mean-aggregation methods by 7\% AUC.
\end{remark}

\subsubsection{Method Selection Theorem}

\begin{theorem}[Optimal Method Selection]
\label{thm:method_selection}
Given a graph $\mathcal{G}$ with diagnostic metrics $(\rho_{\text{FS}}, \delta_{\text{agg}}, h)$, the optimal GNN method selection follows:

\begin{enumerate}
    \item \textbf{High Dilution Regime} ($\delta_{\text{agg}} > \tau_\delta$): \\
    Use sampling (GraphSAGE) or concatenation (H2GCN, MixHop) methods.
    \begin{itemize}
        \item \textit{Rationale}: By Theorems~\ref{thm:sampling} and \ref{thm:concat}, these methods bound or eliminate dilution.
    \end{itemize}

    \item \textbf{Low Dilution + High Alignment} ($\delta_{\text{agg}} \leq \tau_\delta$ and $\rho_{\text{FS}} > \tau_\rho$): \\
    Use feature-aware attention (NAA).
    \begin{itemize}
        \item \textit{Rationale}: Rich feature information can be leveraged without severe dilution.
    \end{itemize}

    \item \textbf{Negative Alignment / Heterophily} ($\rho_{\text{FS}} < 0$ or $h < 0.5$): \\
    Use heterophily-aware methods (H2GCN with 2-hop aggregation).
    \begin{itemize}
        \item \textit{Rationale}: 1-hop neighbors are adversarial; 2-hop neighbors may be same-class.
    \end{itemize}

    \item \textbf{Neutral Regime} ($\delta_{\text{agg}} \leq \tau_\delta$, $|\rho_{\text{FS}}| < \tau_\rho$, $h \geq 0.5$): \\
    Standard GCN/GAT suffices.
    \begin{itemize}
        \item \textit{Rationale}: No pathological conditions; simpler methods are efficient.
    \end{itemize}
\end{enumerate}

Recommended thresholds based on empirical validation: $\tau_\delta \approx 5$, $\tau_\rho \approx 0.1$.
\end{theorem}

\subsubsection{Information-Theoretic Justification for Threshold Selection}
\label{sec:threshold_justification}

We provide rigorous information-theoretic justification for the dilution threshold $\tau_\delta = 5$ (and the more conservative $\tau_\delta = 10$ used in practice).

\begin{theorem}[Information-Theoretic Dilution Bound]
\label{thm:info_dilution}
Let $\mathbf{h}_i$ be a node's ego representation and $\bar{\mathbf{h}}_{\mathcal{N}(i)}$ be the mean neighbor representation after one layer of GCN aggregation. Define the effective signal-to-noise ratio (SNR) as:
\begin{equation}
\text{SNR}_i = \frac{\|\text{proj}_{\mathbf{h}_i^*}(\mathbf{h}_i^{(1)})\|^2}{\|\mathbf{h}_i^{(1)} - \text{proj}_{\mathbf{h}_i^*}(\mathbf{h}_i^{(1)})\|^2}
\end{equation}
where $\mathbf{h}_i^*$ is the optimal representation for classification and $\mathbf{h}_i^{(1)}$ is the post-aggregation representation.

Then, for GCN-style mean aggregation:
\begin{equation}
\text{SNR}_i \leq \frac{1}{d_i \cdot (1 - S_{\cos}(\mathbf{h}_i, \bar{\mathbf{h}}_{\mathcal{N}(i)}))} = \frac{1}{\delta_{\text{agg}}(i)}
\end{equation}
\end{theorem}

\begin{proof}
The GCN output decomposes as $\mathbf{h}_i^{(1)} = \frac{1}{d_i+1}\mathbf{h}_i\mathbf{W} + \frac{d_i}{d_i+1}\bar{\mathbf{h}}_{\mathcal{N}(i)}\mathbf{W}$. For high-degree nodes ($d_i \gg 1$), the ego contribution is $\approx \frac{1}{d_i}\mathbf{h}_i\mathbf{W}$ while the neighbor contribution dominates.

Assuming the ego signal $\mathbf{h}_i$ is aligned with the optimal representation $\mathbf{h}_i^*$, and neighbor signals contribute noise proportional to $(1 - S_{\cos}(\mathbf{h}_i, \bar{\mathbf{h}}_{\mathcal{N}(i)}))$, the SNR bound follows from the signal-noise decomposition.
\end{proof}

\begin{corollary}[Critical Dilution Threshold]
\label{cor:critical_threshold}
When $\delta_{\text{agg}} > 10$, the effective SNR falls below 0.1 (equivalently, -10 dB), meaning:
\begin{itemize}
    \item More than 90\% of the post-aggregation representation is noise from neighbors
    \item The ego signal is effectively ``drowned out'' by aggregation
    \item Classification must rely primarily on neighbor patterns rather than ego features
\end{itemize}

This provides the information-theoretic justification for the threshold $\tau_\delta \in [5, 10]$:
\begin{itemize}
    \item $\tau_\delta = 5$: Conservative threshold (SNR $< 0.2$, 83\% noise)
    \item $\tau_\delta = 10$: Critical threshold (SNR $< 0.1$, 90\% noise)
\end{itemize}
\end{corollary}

\textbf{Empirical Validation}: Our datasets exhibit a clear phase transition around $\delta_{\text{agg}} \approx 10$:
\begin{itemize}
    \item Elliptic ($\delta_{\text{agg}} = 0.9$): NAA improves AUC by 23\%
    \item IEEE-CIS ($\delta_{\text{agg}} = 11.2$): H2GCN outperforms NAA by 7\%
    \item YelpChi ($\delta_{\text{agg}} = 12.6$): H2GCN matches or exceeds NAA
\end{itemize}

The transition at $\delta_{\text{agg}} \approx 10$ is not arbitrary---it corresponds to the information-theoretic boundary where ego signal becomes unrecoverable under mean aggregation.

\textbf{Validation}: Table~\ref{tab:method_selection_validation} shows that this theorem correctly predicts the best method on all five experimental datasets.

\begin{table}[h]
\centering
\caption{Method Selection Theorem Validation}
\label{tab:method_selection_validation}
\small
\begin{tabular}{lccccc}
\toprule
\textbf{Dataset} & $\delta_{\text{agg}}$ & $\rho_{\text{FS}}$ & \textbf{Regime} & \textbf{Predicted} & \textbf{Actual} \\
\midrule
Elliptic & 0.9 & 0.28 & Low-$\delta$, High-$\rho$ & NAA & NAA \\
Amazon & -- & 0.18 & High-$\rho$ & NAA & NAA \\
YelpChi & 12.6 & 0.01 & High-$\delta$ & SAGE/H2GCN & H2GCN \\
IEEE-CIS & 11.2 & 0.06 & High-$\delta$ & SAGE/H2GCN & H2GCN \\
DGraphFin & -- & -- & Extreme imbalance & None & None \\
\bottomrule
\end{tabular}
\end{table}

\subsection{Unified Aggregation Framework}

Based on the two similarity notions, we propose a unified aggregation framework.

\begin{definition}[Feature-Structure Disentangled Attention]
\label{def:fsd_attention}
The FSD attention mechanism computes edge weights as:
\begin{equation}
\alpha_{ij} = \text{softmax}_j\left(e_{ij}^{\text{struct}} + \lambda \cdot e_{ij}^{\text{feat}}\right)
\end{equation}
where:
\begin{itemize}
    \item $e_{ij}^{\text{struct}}$: Structure-based attention score (e.g., learned attention in GAT)
    \item $e_{ij}^{\text{feat}}$: Feature-based attention score (e.g., numerical similarity in NAA)
    \item $\lambda \geq 0$: Trade-off parameter controlling the balance
\end{itemize}
\end{definition}

\textbf{Existing Methods as Special Cases:}

\begin{table}[h]
\centering
\caption{Existing GNN methods as special cases of the FSD framework}
\label{tab:fsd_special_cases}
\small
\begin{tabular}{lccc}
\toprule
\textbf{Method} & $e_{ij}^{\text{struct}}$ & $e_{ij}^{\text{feat}}$ & $\lambda$ \\
\midrule
GCN & $1/\sqrt{d_i d_j}$ & -- & 0 \\
GAT & $\mathbf{a}^\top[\mathbf{W}\mathbf{h}_i \| \mathbf{W}\mathbf{h}_j]$ & -- & 0 \\
GATv2 & $\mathbf{a}^\top \text{LeakyReLU}(\mathbf{W}[\mathbf{h}_i \| \mathbf{h}_j])$ & -- & 0 \\
\textbf{NAA (Ours)} & Learned attention & NumericalSim$(\mathbf{x}_i, \mathbf{x}_j)$ & $\beta$ \\
H2GCN & Multi-hop separation & Implicit via 2-hop & Implicit \\
\bottomrule
\end{tabular}
\end{table}

This unification reveals that:
\begin{enumerate}
    \item \textbf{Standard GNNs} (GCN, GAT, GATv2) use only structural information ($\lambda = 0$)
    \item \textbf{NAA} explicitly incorporates feature similarity with learned $\lambda = \beta$
    \item \textbf{H2GCN} achieves implicit feature-awareness through multi-hop separation, which allows aggregating from 2-hop neighbors that may be more feature-similar under heterophily
\end{enumerate}

\subsection{Theoretical Analysis: When Feature-Attention Helps}

We now analyze when incorporating feature-based attention ($\lambda > 0$) improves over pure structural attention ($\lambda = 0$).

\begin{definition}[Information Content]
For node classification with labels $Y$:
\begin{itemize}
    \item \textbf{Feature Information}: $I_{\text{feat}} = I(\mathbf{X}; Y)$ --- mutual information between features and labels
    \item \textbf{Structure Information}: $I_{\text{struct}} = I(\mathcal{A}; Y)$ --- mutual information between adjacency and labels
\end{itemize}
\end{definition}

\begin{theorem}[Optimal Aggregation Regime]
\label{thm:optimal_regime}
Consider binary node classification on graph $\mathcal{G}$ with features $\mathbf{X}$ and labels $Y$. Let $\rho_{\text{FS}}$ be the Feature-Structure Alignment Score. The optimal aggregation strategy depends on three regimes:

\textbf{Regime A (Feature-Dominant):} When $I_{\text{feat}} / I_{\text{struct}} > \tau_1$ and $\rho_{\text{FS}} < \tau_2$:
\begin{itemize}
    \item Feature-weighted attention ($\lambda > 0$) outperforms pure structural attention
    \item \textbf{Recommended}: NAA-style methods
\end{itemize}

\textbf{Regime B (Structure-Dominant):} When $I_{\text{struct}} / I_{\text{feat}} > \tau_1$ and $\rho_{\text{FS}} > \tau_2$:
\begin{itemize}
    \item Structural attention is sufficient
    \item \textbf{Recommended}: Standard GCN/GAT
\end{itemize}

\textbf{Regime C (Heterophily):} When $\rho_{\text{FS}} < 0$ (structural neighbors are feature-dissimilar):
\begin{itemize}
    \item Explicit disentanglement is necessary
    \item \textbf{Recommended}: H2GCN-style multi-hop separation or high $\lambda$
\end{itemize}
\end{theorem}

\begin{proof}
We analyze the expected prediction error under different aggregation strategies.

\textbf{Setup}: Consider a simplified linear model where the aggregated representation is:
\begin{equation}
\tilde{\mathbf{h}}_i = (1-\alpha) \mathbf{h}_i + \alpha \sum_{j \in \mathcal{N}(i)} w_{ij} \mathbf{h}_j
\end{equation}
where $\alpha \in [0,1]$ controls aggregation strength and $w_{ij}$ are attention weights summing to 1.

\textbf{Aggregation Error Analysis}: For node classification, the prediction depends on how well the aggregated representation preserves label-relevant information. Define the aggregation-induced noise for node $i$:
\begin{equation}
\mathcal{L}_{\text{agg}}(i) = \mathbb{E}\left[\left\| \tilde{\mathbf{h}}_i - \mathbf{h}_i^* \right\|^2\right]
\end{equation}
where $\mathbf{h}_i^*$ is the optimal representation for predicting $y_i$.

\textbf{Case 1 (Pure Structural Aggregation, $\lambda = 0$)}: Attention weights are based only on structure:
\begin{equation}
w_{ij}^{\text{struct}} = \frac{\exp(e_{ij}^{\text{struct}})}{\sum_{k \in \mathcal{N}(i)} \exp(e_{ik}^{\text{struct}})}
\end{equation}

The expected noise from structural neighbors is:
\begin{align}
\mathcal{L}_{\text{struct}}(i) &= \alpha^2 \mathbb{E}_{j \sim w^{\text{struct}}}\left[\|\mathbf{h}_j - \mathbf{h}_i^*\|^2\right] \\
&= \alpha^2 \left( \underbrace{\mathbb{E}[\|\mathbf{h}_j - \mathbf{h}_i\|^2]}_{\text{neighbor variance}} + \underbrace{\|\mathbf{h}_i - \mathbf{h}_i^*\|^2}_{\text{ego error}} \right)
\end{align}

When $\rho_{\text{FS}}$ is low, structural neighbors have dissimilar features, so $\mathbb{E}[\|\mathbf{h}_j - \mathbf{h}_i\|^2]$ is large.

\textbf{Case 2 (Feature-Weighted Aggregation, $\lambda > 0$)}: Attention incorporates feature similarity:
\begin{equation}
w_{ij}^{\text{FSD}} = \frac{\exp(e_{ij}^{\text{struct}} + \lambda \cdot e_{ij}^{\text{feat}})}{\sum_{k \in \mathcal{N}(i)} \exp(e_{ik}^{\text{struct}} + \lambda \cdot e_{ik}^{\text{feat}})}
\end{equation}

The feature similarity term down-weights dissimilar neighbors:
\begin{equation}
\mathcal{L}_{\text{FSD}}(i) = \alpha^2 \mathbb{E}_{j \sim w^{\text{FSD}}}\left[\|\mathbf{h}_j - \mathbf{h}_i^*\|^2\right]
\end{equation}

Since $e_{ij}^{\text{feat}}$ is higher for feature-similar nodes, $w^{\text{FSD}}$ assigns more weight to neighbors with similar features, reducing the variance term.

\textbf{Regime Analysis}:

\textit{Regime A}: When $I_{\text{feat}} > I_{\text{struct}}$ and $\rho_{\text{FS}} < \tau_2$:
\begin{itemize}
\item Labels are primarily determined by features ($\mathbf{h}_i^* \approx f(\mathbf{x}_i)$)
\item Structural neighbors have dissimilar features (low $\rho_{\text{FS}}$)
\item $\mathcal{L}_{\text{struct}} > \mathcal{L}_{\text{FSD}}$ because feature-weighting reduces neighbor variance
\end{itemize}

\textit{Regime B}: When $I_{\text{struct}} > I_{\text{feat}}$ and $\rho_{\text{FS}} > \tau_2$:
\begin{itemize}
\item Graph structure carries label information
\item High $\rho_{\text{FS}}$ means structural neighbors are also feature-similar
\item $\mathcal{L}_{\text{struct}} \approx \mathcal{L}_{\text{FSD}}$, so the simpler structural approach suffices
\end{itemize}

\textit{Regime C}: When $\rho_{\text{FS}} < 0$:
\begin{itemize}
\item Structural neighbors are \textit{adversarially} dissimilar
\item $\mathbb{E}_{(i,j) \in \mathcal{E}}[S_{\text{feat}}] < \mathbb{E}_{(i,j) \notin \mathcal{E}}[S_{\text{feat}}]$
\item Aggregation from 1-hop neighbors actively harms prediction
\item Multi-hop methods (H2GCN) can reach same-label nodes at distance 2
\end{itemize}

\textbf{Threshold Characterization}: The thresholds $\tau_1, \tau_2$ depend on:
\begin{equation}
\tau_1 = \frac{\text{Var}(\mathbf{h}_i | y_i)}{\text{Var}(\mathbf{h}_i | \mathcal{N}(i))}, \quad \tau_2 = \frac{\sigma_{\text{feat}}^2}{\sigma_{\text{struct}}^2}
\end{equation}
where $\sigma_{\text{feat}}^2$ and $\sigma_{\text{struct}}^2$ are the within-class variances on feature and structural graphs respectively.
\end{proof}

\begin{corollary}[NAA's Optimal Conditions]
\label{cor:naa_conditions}
NAA provides significant improvement when:
\begin{enumerate}
    \item \textbf{High-dimensional numerical features} ($d > 100$): Large feature space provides rich information for similarity computation
    \item \textbf{Multi-scale numerical values}: Log-normalization preserves relative differences across magnitudes
    \item \textbf{Moderate Feature-Structure alignment} ($0 < \rho_{\text{FS}} < 0.5$): Features are informative but not fully captured by structure
    \item \textbf{Moderate class imbalance}: Extreme imbalance ($<1\%$ positive) dominates all architectural choices
\end{enumerate}
\end{corollary}

\subsection{Spectral Interpretation}
\label{sec:spectral}

We provide an alternative theoretical perspective using spectral graph theory.

\textbf{Graph Signal Processing View}: Node labels $\mathbf{y} \in \mathbb{R}^N$ can be viewed as a signal on the graph. GNN aggregation acts as a graph filter.

For structural graph $\mathcal{G}_{\text{struct}}$ with normalized Laplacian $\mathbf{L}_{\text{struct}}$, the Graph Fourier Transform of $\mathbf{y}$ is:
\begin{equation}
\hat{\mathbf{y}}_{\text{struct}} = \mathbf{U}_{\text{struct}}^\top \mathbf{y}
\end{equation}
where $\mathbf{U}_{\text{struct}}$ contains eigenvectors of $\mathbf{L}_{\text{struct}}$.

Similarly, we can construct a \textbf{feature similarity graph} $\mathcal{G}_{\text{feat}}$ with edges weighted by $S_{\text{feat}}(i,j)$, yielding $\hat{\mathbf{y}}_{\text{feat}}$.

\begin{proposition}[Spectral Characterization]
\label{prop:spectral}
Let $\lambda_{\text{struct}}$ and $\lambda_{\text{feat}}$ denote the average frequency of label signal $\mathbf{y}$ on $\mathcal{G}_{\text{struct}}$ and $\mathcal{G}_{\text{feat}}$ respectively:
\begin{equation}
\bar{\lambda}_{\text{struct}} = \frac{\sum_k \lambda_k |\hat{y}_k^{\text{struct}}|^2}{\sum_k |\hat{y}_k^{\text{struct}}|^2}
\end{equation}

\begin{itemize}
    \item If $\bar{\lambda}_{\text{struct}} > \bar{\lambda}_{\text{feat}}$: Labels are smoother on $\mathcal{G}_{\text{feat}}$ --- feature-based aggregation is preferable
    \item If $\bar{\lambda}_{\text{struct}} < \bar{\lambda}_{\text{feat}}$: Labels are smoother on $\mathcal{G}_{\text{struct}}$ --- structural aggregation is preferable
\end{itemize}
\end{proposition}

\textbf{Intuition}: GNN aggregation is a low-pass filter. Low-pass filtering is effective when the target signal is low-frequency (smooth) on the graph. If labels are high-frequency on $\mathcal{G}_{\text{struct}}$ (heterophily) but low-frequency on $\mathcal{G}_{\text{feat}}$ (feature-label alignment), aggregating on $\mathcal{G}_{\text{feat}}$ is more appropriate.

\subsection{Empirical Predictions}

Our theoretical framework makes concrete predictions about when different methods should excel. Table~\ref{tab:predictions} shows predictions based on the extended FSD framework with three diagnostic metrics.

\begin{table}[h]
\centering
\caption{Theoretical predictions for our experimental datasets based on the extended FSD framework}
\label{tab:predictions}
\small
\begin{tabular}{lccccp{2.8cm}}
\toprule
\textbf{Dataset} & \textbf{Dim} & \textbf{$\rho_{\text{FS}}$} & \textbf{$\delta_{\text{agg}}$} & \textbf{Predicted Winner} & \textbf{Reasoning} \\
\midrule
Elliptic & 165 & 0.28 & 0.9 (Low) & NAA & High-dim, low dilution \\
Amazon & 767 & 0.18 & -- & NAA & Very high-dim \\
YelpChi & 32 & 0.01 & 12.6 (High) & H2GCN/SAGE & Low-dim, high dilution \\
DGraphFin & 17 & N/A & N/A & None & Extreme imbalance \\
\midrule
\textbf{IEEE-CIS} & 394 & 0.06 & 11.2 (High) & \textbf{SAGE/H2GCN} & High dilution dominates \\
\bottomrule
\end{tabular}
\end{table}

\textbf{Key Insight}: The IEEE-CIS dataset serves as a \textit{critical test case} for the extended framework. Based on $\rho_{\text{FS}} = 0.06$ alone, we would predict ``no clear winner'' (similar to the mixed regime). However, incorporating $\delta_{\text{agg}} = 11.2$ reveals high aggregation dilution, predicting that sampling/concatenation methods (GraphSAGE, H2GCN) should outperform mean-aggregation methods (GCN, GAT, NAA).

We validate these predictions empirically in Section~\ref{sec:experiments}.
